\documentclass[11pt,a4paper]{article}
\usepackage[T1]{fontenc}
\usepackage{lmodern,microtype}
\usepackage[margin=25mm]{geometry}
\usepackage{amsmath,amssymb,amsthm,mathtools}
\usepackage{booktabs,array,enumitem,xcolor,fancyhdr,needspace}
\usepackage[hypertexnames=false,colorlinks=true,linkcolor=blue!45!black,citecolor=blue!45!black,urlcolor=blue!45!black]{hyperref}
\setlength{\emergencystretch}{3em}
\setlength{\parskip}{0.25em}
\setlist{topsep=3pt,itemsep=3pt}
\pagestyle{fancy}\fancyhf{}
\fancyhead[L]{\small Complete monotonicity of arctangent inverses}
\fancyhead[R]{\small Research manuscript}
\fancyfoot[C]{\thepage}
\setlength{\headheight}{14pt}
\newtheorem{theorem}{Theorem}[section]
\newtheorem{lemma}[theorem]{Lemma}
\newtheorem{proposition}[theorem]{Proposition}
\newtheorem{corollary}[theorem]{Corollary}
\theoremstyle{definition}\newtheorem{definition}[theorem]{Definition}
\theoremstyle{remark}\newtheorem{remark}[theorem]{Remark}
\newcommand{\CM}{\mathrm{CM}}
\newcommand{\BF}{\mathrm{BF}}
\newcommand{\Sclass}{\mathcal S}
\newcommand{\dd}{\,\mathrm d}
\DeclareMathOperator{\artanh}{artanh}
\title{Complete Monotonicity of the Inverses\\of Two Arctangent Quotients\\[5pt]
\large A sharp parameter classification and a positive cosine representation}
\author{Keio Research Notes}
\date{31 October 2026}
\hypersetup{pdftitle={Complete Monotonicity of the Inverses of Two Arctangent Quotients},pdfauthor={Keio Research Notes},pdfsubject={Explicit proofs addressing Question 1 of Keady, Khajohnsaksumeth, and Wiwatanapataphee (2018)}}
\begin{document}
\maketitle
\begin{abstract}
Let $\mu$ be the positive functional inverse of
$x\mapsto \arctan(1/x)/x$. We prove that $1/\mu$ is a Bernstein
function and that $\mu^p$ is strictly completely monotone for every
$p>0$. Together with a short moment-rigidity argument showing that
$\mu$ is not Stieltjes, this gives complete mathematical answers to
all clauses of Question 1 of Keady, Khajohnsaksumeth, and
Wiwatanapataphee (2018). The main calculation classifies the family
$U_\lambda(x)=[\arctan x+\lambda x/(1+x^2)]^{-1}$ for
$\lambda\ge0$: it is completely monotone exactly when
$0\le\lambda\le1$. An explicit Stieltjes density in the squared
variable is compared with the square-root resolvent. An exact mass
identity then converts the resulting signed cosine transform into
a nonnegative $1-\cos$ kernel. An existing inverse differential-equation
criterion transfers this positivity to the original inverses.
Every assertion about derivative signs is proved for all orders;
no numerical sign test or unproved conjecture enters the proof.
Publication priority is discussed separately from mathematical validity.
\end{abstract}

\section{The original question and the result}

For $z>0$, set
\begin{equation}\label{eq:original}
 \phi_1(z)=\frac{\arctan(1/z)}{z},\qquad
 \phi_2(z)=\frac{\arctan(1/\sqrt z)}{\sqrt z},\qquad
 \mu=\phi_1^{-1}.
\end{equation}
Both maps are strictly decreasing bijections of $(0,\infty)$ onto
itself. Also $\phi_2(z)=\phi_1(\sqrt z)$, so
\begin{equation}\label{eq:inverse-square}
 \phi_2^{-1}(y)=\mu(y)^2.
\end{equation}
Question 1 on printed page 2 of \cite{KKW2018} asks whether each
of these two inverse functions is completely monotone, and whether
$\mu$ might even be Stieltjes. These are the original functions,
domains, and quantifiers of that question.

\begin{definition}
A smooth function $f:(0,\infty)\to\mathbb R$ is completely monotone
if $(-1)^n f^{(n)}(x)\ge0$ for every integer $n\ge0$ and every $x>0$.
A nonnegative smooth function $h$ is Bernstein if $h'$ is completely
monotone. A Stieltjes function has a representation
\[
 f(x)=b+\int_{[0,\infty)}\frac{\rho(\dd t)}{x+t},
 \quad b\ge0,\quad \rho\ge0,\quad
 \int_{[0,\infty)}\frac{\rho(\dd t)}{1+t}<\infty.
\]
An atom at zero includes the usual $a/x$ term.
\end{definition}
These standard classes and their integral representations are treated
systematically in \cite{SSV2012}.

\Needspace{15\baselineskip}
\begin{theorem}[Answers to the original question]\label{thm:main}
For the function $\mu$ in \eqref{eq:original}:
\begin{enumerate}[label=\textup{(\roman*)}]
\item $h=1/\mu$ is a Bernstein function. More precisely,
$(-1)^{n-1}h^{(n)}(y)>0$ for all $n\ge1$ and $y>0$.
\item For every real $p>0$,
\[
 (-1)^n\frac{\mathrm d^n}{\mathrm dy^n}\mu(y)^p>0
 \qquad(n\ge0,\ y>0).
\]
In particular, both $\phi_1^{-1}$ and $\phi_2^{-1}$ are completely
monotone. The function $\mu$ is also logarithmically completely
monotone.
\item $\mu$ is not a Stieltjes function. Neither is $\mu^2$.
\end{enumerate}
\end{theorem}

The key change of variable is exact. From
$y\mu(y)=\arctan(1/\mu(y))$ one obtains
\begin{equation}\label{eq:h-inverse}
 y=h(y)\arctan h(y).
\end{equation}
Thus $h$ is the increasing functional inverse of $g(x)=x\arctan x$.
Its differential equation is
\begin{equation}\label{eq:h-ode}
 h'(y)=U_1(h(y)),\qquad
 U_1(x)=\frac1{\arctan x+x/(1+x^2)}.
\end{equation}
The central task is therefore an explicit positivity theorem for
$U_1$, not an assumption that inverses of arbitrary completely
monotone functions preserve that class.

The argument below proves a sharp family result before applying it
to \eqref{eq:h-ode}. The inverse-ODE implication itself is an existing
criterion; we include its elementary sign proof in Section~\ref{sec:ode}.
Related reciprocal-arctangent work is discussed in
Section~\ref{sec:literature}. The mathematical conclusions here do
not by themselves establish publication priority.

\section{An auxiliary classification problem}

Put $L=\pi/2$. For a parameter $\lambda\ge0$, define
\begin{align}
 F_\lambda(x)&=\arctan x+\frac{\lambda x}{1+x^2},
 &U_\lambda(x)&=\frac1{F_\lambda(x)},\label{eq:F-U}\\
 g_\lambda(x)&=x\arctan x+
           \frac{\lambda-1}{2}\log(1+x^2).
 &&\label{eq:g-lambda}
\end{align}
The auxiliary problem is to determine all $\lambda\ge0$ for which
$U_\lambda$ is completely monotone. The endpoint $\lambda=1$ is
precisely the function required by the original problem.

\begin{theorem}[Sharp classification]\label{thm:classification}
For $\lambda\ge0$ the following are equivalent:
\begin{enumerate}[label=\textup{(\roman*)}]
\item $U_\lambda$ is completely monotone on $(0,\infty)$;
\item the increasing inverse $h_\lambda=g_\lambda^{-1}$ is a
Bernstein function on $(0,\infty)$;
\item $0\le\lambda\le1$.
\end{enumerate}
In this range all alternating derivative signs of $U_\lambda$ and
$h_\lambda'$ are strict. Every $h_\lambda^{-p}$, $p>0$, is strictly
completely monotone.
\end{theorem}

The implication from (i) to (ii) is not asserted to be reversible
for general functions. Necessity for this particular family will
be proved directly from $g_\lambda''$.

\section{Converting a spectral deficit into a positive kernel}

The following elementary comparison is useful because the cosine
transform in its hypothesis need not have a sign.

\begin{lemma}[Cosine compensation]\label{lem:compensation}
Suppose $D\ge0$ belongs to $L^1(0,\infty)$ and
\begin{equation}\label{eq:comparison}
 U(x)=c+\frac a x-\frac2\pi\int_0^\infty
              D(r)\frac{x}{x^2+r^2}\dd r,
 \qquad x>0,
\end{equation}
where $c,a\ge0$. Then
\begin{equation}\label{eq:mass-coefficient}
 A:=\lim_{x\to\infty}x\bigl(U(x)-c\bigr)
   =a-\frac2\pi\int_0^\infty D(r)\dd r.
\end{equation}
If $A\ge0$, then $U$ is completely monotone, with representation
\begin{equation}\label{eq:positive-laplace}
 U(x)=c+\int_0^\infty e^{-xt}k(t)\dd t,\qquad
 k(t)=A+\frac2\pi\int_0^\infty D(r)(1-\cos rt)\dd r\ge0.
\end{equation}
If $A<0$, then $U$ is not completely monotone.
\end{lemma}

\begin{proof}
Multiplying \eqref{eq:comparison} by $x$ after subtracting $c$,
and using $x^2/(x^2+r^2)\uparrow1$, proves
\eqref{eq:mass-coefficient} by monotone convergence. Substitution
of that identity into \eqref{eq:comparison} gives
\[
 U(x)=c+\frac A x+\frac2\pi\int_0^\infty D(r)
                  \left(\frac1x-\frac{x}{x^2+r^2}\right)\dd r.
\]
For $x>0$ and $r\ge0$,
\begin{equation}\label{eq:cosine-identity}
 \frac1x-\frac{x}{x^2+r^2}
 =\int_0^\infty e^{-xt}(1-\cos rt)\dd t.
\end{equation}
Tonelli's theorem now yields \eqref{eq:positive-laplace} when
$A\ge0$. The density is bounded by
$A+(4/\pi)\|D\|_1$. Consequently all differentiations under its
Laplace integral are justified on compact subintervals of $x>0$,
and give the required signs for every order. If $A<0$, then
$U(x)<c$ for all sufficiently large $x$, whereas $U(x)\to c$.
A decreasing function with limit $c$ cannot have this property.
\end{proof}

\section{The explicit spectral calculation}

For $r>1$ set
\begin{equation}\label{eq:b}
 b_\lambda(r)=\artanh(1/r)-\frac{\lambda r}{r^2-1},
 \qquad
 \artanh(1/r)=\frac12\log\frac{r+1}{r-1}.
\end{equation}
Define a nonnegative function on $(0,\infty)$ by
\begin{equation}\label{eq:D}
 D_\lambda(r)=
 \begin{cases}
 L^{-1},&0<r\le1,\\[2pt]
 \displaystyle\frac{b_\lambda(r)^2}
 {L\bigl(L^2+b_\lambda(r)^2\bigr)},&r>1.
 \end{cases}
\end{equation}
The two pieces meet continuously at $r=1$. Moreover,
\[
 b_\lambda(r)=\frac{1-\lambda}{r}
       +\frac{1/3-\lambda}{r^3}+O(r^{-5})
       \qquad(r\to\infty),
\]
so $D_\lambda(r)=O(r^{-2})$ and $D_\lambda\in L^1(0,\infty)$.
For $\lambda=1$ the stronger estimate is $D_1(r)=O(r^{-6})$.

\begin{lemma}[Stieltjes representation in the squared variable]
\label{lem:K}
For $\lambda\ge0$ and $w\in\mathbb C\setminus(-\infty,0]$, using
the principal square root, let
\[
 K_\lambda(w)=
 \frac1{\sqrt w\,\arctan\sqrt w+\lambda w/(1+w)}.
\]
Then, with $a_\lambda=1/(1+\lambda)$,
\begin{equation}\label{eq:K-representation}
 K_\lambda(w)=\frac{a_\lambda}{w}
 +\frac{2L}{\pi}\int_1^\infty
    \frac{\dd r}{(w+r^2)\bigl(L^2+b_\lambda(r)^2\bigr)}.
\end{equation}
\end{lemma}

\begin{proof}
Write the denominator as
\begin{equation}\label{eq:Q}
 Q_\lambda(w)=\int_0^1\frac{w}{1+wt^2}\dd t
                     +\frac{\lambda w}{1+w}.
\end{equation}
For $\operatorname{Im}w>0$,
\[
 \operatorname{Im}\frac{w}{1+wt^2}
   =\frac{\operatorname{Im}w}{|1+wt^2|^2}>0.
\]
The integral also gives an analytic continuation to
$\mathbb C\setminus(-\infty,-1]$. Thus $Q_\lambda$ has no zero in the upper half-plane; conjugation
gives the corresponding statement below. It is positive on the
positive real axis. For $-1<w<0$ the integral in
\eqref{eq:Q} is negative and the remaining term is nonpositive,
so there is no zero there either.

At zero,
\[
 Q_\lambda(w)=(1+\lambda)w+O(w^2),\qquad
 K_\lambda(w)=a_\lambda/w+O(1).
\]
At $w=-1$, $K_\lambda$ is bounded and tends to zero. When
$\lambda>0$ this follows from the simple pole of
$\lambda w/(1+w)$ dominating the logarithmic term. When
$\lambda=0$ the logarithmic divergence itself gives the claim.
At infinity,
\[
 K_\lambda(w)=O(|w|^{-1/2})
\]
uniformly in the slit plane. Indeed, if $z=\sqrt w$ then
$\operatorname{Re}z>0$ and
$\arctan z=L-\arctan(1/z)$, so the convergent expansion in
$1/z$ applies uniformly for large $|z|$.

For the upper boundary point $w=-r^2+i0$ with $r>1$,
\[
 \arctan(ir+0)=L+i\artanh(1/r),\qquad
 F_\lambda(ir+0)=L+i b_\lambda(r).
\]
Consequently
\begin{equation}\label{eq:boundary-K}
 K_\lambda(-r^2+i0)
 =\frac1{-r b_\lambda(r)+iLr},\qquad
 -\frac1\pi\operatorname{Im}K_\lambda(-r^2+i0)
 =\frac{L}{\pi r\bigl(L^2+b_\lambda(r)^2\bigr)}.
\end{equation}
On $-1<w<0$ the boundary value is real.

Apply Cauchy's integral formula to
$K_\lambda(w)-a_\lambda/w$ on a keyhole contour around the
negative axis. The large circle contributes $O(R^{-1/2})$.
The circle at zero contributes zero in the limit after removal of
the pole, and a small indentation at $-1$ contributes zero by
boundedness. The two sides of the cut are conjugate, so their
difference gives the density in \eqref{eq:boundary-K}.
It is locally integrable at $r=1$ and gives an integrand
$O(r^{-2})$ at infinity. Changing the cut variable from $s=r^2$
to $r$ yields \eqref{eq:K-representation}. This contour argument
also excludes any omitted atom, pole, or additive term.
\end{proof}

\Needspace{12\baselineskip}
\begin{proposition}[Positive Laplace density and exact mass]
\label{prop:kernel}
For every $\lambda\ge0$,
\begin{equation}\label{eq:mass}
 \frac2\pi\int_0^\infty D_\lambda(r)\dd r
 =\frac1{1+\lambda}-\frac{1-\lambda}{L^2}.
\end{equation}
For $0\le\lambda\le1$,
\begin{align}
 U_\lambda(x)&=\frac1L+
           \int_0^\infty e^{-xt}k_\lambda(t)\dd t,
          \label{eq:U-density}\\
 k_\lambda(t)&=\frac{1-\lambda}{L^2}
       +\frac2\pi\int_0^\infty
                     D_\lambda(r)(1-\cos rt)\dd r\ge0.
          \label{eq:k-density}
\end{align}
In fact $k_\lambda(t)>0$ for every $t>0$.
\end{proposition}

\begin{proof}
Since $U_\lambda(x)=xK_\lambda(x^2)$, Lemma~\ref{lem:K} gives
\[
 U_\lambda(x)=\frac{a_\lambda}{x}
 +\frac{2L}{\pi}\int_1^\infty
 \frac{x\dd r}{(x^2+r^2)(L^2+b_\lambda(r)^2)}.
\]
Subtract the elementary reference integral
\[
 \frac1L=\frac2{\pi L}\int_0^\infty
                         \frac{x}{x^2+r^2}\dd r.
\]
The result is precisely \eqref{eq:comparison}, with $c=1/L$,
$a=a_\lambda$, and $D=D_\lambda$. On the other hand,
\[
 F_\lambda(x)=L+\frac{\lambda-1}{x}+O(x^{-3}),\qquad
 U_\lambda(x)=\frac1L+\frac{1-\lambda}{L^2x}+O(x^{-2}).
\]
Lemma~\ref{lem:compensation} gives \eqref{eq:mass} and
\eqref{eq:U-density}--\eqref{eq:k-density}.
For $t>0$, the continuous nonnegative function $1-\cos rt$ is
not identically zero on $0<r<1$, where $D_\lambda=1/L$.
Thus $k_\lambda(t)>0$.
\end{proof}

It follows at once that $U_\lambda$ is strictly completely monotone
for $0\le\lambda\le1$. If $\lambda>1$, the displayed expansion
shows $U_\lambda(x)<1/L$ for all sufficiently large $x$, although
its limit is $1/L$. Hence it cannot be completely monotone. This
proves the equivalence of (i) and (iii) in
Theorem~\ref{thm:classification}.

At the endpoint needed for the original problem, the mass identity
and the kernel take the especially simple form
\begin{equation}\label{eq:endpoint}
 \int_0^\infty D_1(r)\dd r=\frac\pi4,\qquad
 U_1(x)=\frac2\pi+\frac2\pi\int_0^\infty e^{-xt}
       \left[\int_0^\infty D_1(r)(1-\cos rt)\dd r\right]\dd t.
\end{equation}

\section{Passing to the functional inverse}\label{sec:ode}

We record the existing inverse-ODE principle in the exact form
needed here. The proof is included to make every derivative-order
claim independent of numerical evidence. See also Li's
inverse-ODE criterion \cite[Lemma~2.6, version~2]{Li2026} and its use in
\cite[APP-0066]{Pudim2026}.

\begin{lemma}[Inverse-ODE sign induction]\label{lem:ode}
Let $U$ be a positive completely monotone function on $(0,\infty)$.
If a positive smooth function $h$ on $(0,\infty)$ satisfies
$h'=U(h)$, then $h$ is Bernstein. If all alternating signs of
$U$ are strict, so are those of $h'$. Furthermore,
$h^{-p}$ is completely monotone for every $p>0$.
\end{lemma}

\begin{proof}
The base case is $h'=U(h)>0$. Suppose
$(-1)^{j-1}h^{(j)}\ge0$ for $1\le j\le n$. The formula for the
$n$th derivative of $U(h)$ is a sum of terms
\begin{equation}\label{eq:faa}
 \frac{n!}{\prod_{j=1}^n m_j!(j!)^{m_j}}
 U^{(k)}(h)\prod_{j=1}^n\bigl(h^{(j)}\bigr)^{m_j},
 \qquad \sum_{j=1}^n j m_j=n,\quad
 k=\sum_{j=1}^n m_j.
\end{equation}
Every coefficient is positive. The sign contributed by
$U^{(k)}$ is $(-1)^k$, and the product contributes
$(-1)^{\sum(j-1)m_j}=(-1)^{n-k}$. Thus every term has sign
$(-1)^n$. Since $(U\circ h)^{(n)}=h^{(n+1)}$, the induction
closes for every $n$. Strictness follows, for example, from the
term $U^{(n)}(h)(h')^n$.

Apply the same formula to $v(h)$ with $v(x)=x^{-p}$, whose
derivatives have strict alternating signs. Using the signs of
$h^{(j)}$ just proved, each term in its $n$th derivative has
sign $(-1)^n$. The term $v^{(n)}(h)(h')^n$ gives strictness.
\end{proof}

\begin{proof}[Completion of Theorem~\ref{thm:classification}]
For every $\lambda\ge0$, $g_\lambda'=F_\lambda>0$,
$g_\lambda(0)=0$, and $g_\lambda(x)\to\infty$. Hence
$h_\lambda$ is a smooth increasing inverse on $(0,\infty)$ and
$h_\lambda'=U_\lambda(h_\lambda)$. When $0\le\lambda\le1$,
Proposition~\ref{prop:kernel} and Lemma~\ref{lem:ode} prove all
the asserted positive conclusions.

For $\lambda>1$, direct differentiation gives
\[
 g_\lambda''(x)
 =\frac{(1+\lambda)+(1-\lambda)x^2}{(1+x^2)^2}<0
 \quad\text{if }x^2>\frac{1+\lambda}{\lambda-1}.
\]
Therefore $h_\lambda''(g_\lambda(x))
=-g_\lambda''(x)/(g_\lambda'(x))^3>0$ there. A Bernstein
function must have nonpositive second derivative. This proves
the remaining necessity without using a converse to
Lemma~\ref{lem:ode}.
\end{proof}

\begin{proof}[Parts (i) and (ii) of Theorem~\ref{thm:main}]
Take $\lambda=1$ in Theorem~\ref{thm:classification} and use
\eqref{eq:h-inverse}. This gives $h=1/\mu$ Bernstein and all
the stated signs for $\mu^p=h^{-p}$, including $p=1,2$.
For logarithmic complete monotonicity, observe that
\[
 -\frac{\mathrm d}{\mathrm dy}\log\mu(y)
 =\frac{h'(y)}{h(y)}.
\]
Both factors $h'$ and $1/h$ are completely monotone. Their product
is completely monotone by Leibniz's formula. Hence
$(-1)^n(\log\mu)^{(n)}\ge0$ for all $n\ge1$.
\end{proof}

\section{The Stieltjes clause}

The following argument was obtained in an earlier local research
note in this project. We include the complete proof so that the
entire original Question 1 is answered within this single manuscript.

\begin{lemma}[Three-coefficient rigidity]\label{lem:rigidity}
If a Stieltjes function $V$ satisfies
\[
 V(y)=\frac c{y+d}+o(y^{-3})\quad(y\to\infty),
 \qquad c>0,\ d\ge0,
\]
then $V(y)=c/(y+d)$ for every $y>0$.
\end{lemma}

\begin{proof}
Since $V(y)\to0$, its Stieltjes representation has zero additive
constant. Write $V(y)=\int(y+t)^{-1}\rho(\dd t)$. Monotone
convergence in $y/(y+t)$ gives total mass $c$. Once this is known,
monotone convergence in $yt/(y+t)$ gives first moment $cd$ from
\[
 y^2\bigl(c/y-V(y)\bigr)
 =\int\frac{yt}{y+t}\rho(\dd t)\longrightarrow cd.
\]
With that first moment established, the same argument gives the
second moment:
\[
 y^3\bigl(V(y)-c/y+cd/y^2\bigr)
 =\int\frac{yt^2}{y+t}\rho(\dd t)\longrightarrow cd^2.
\]
No moment was assumed finite before being justified. Thus
$\int(t-d)^2\rho(\dd t)=0$. Positivity forces
$\rho=c\delta_d$, which proves the assertion.
\end{proof}

\begin{proof}[Part (iii) of Theorem~\ref{thm:main}]
The defining equation gives
\[
 y\mu=L-\arctan\mu,\qquad 0<\mu<L/y.
\]
For $u>0$,
\[
 0<u-\arctan u=\int_0^u\frac{t^2}{1+t^2}\dd t<\frac{u^3}{3}.
\]
It follows that for every $y>0$,
\begin{equation}\label{eq:strict-asymptotic}
 0<\mu(y)-\frac L{y+1}
   =\frac{\mu(y)-\arctan\mu(y)}{y+1}
   <\frac{L^3}{3y^3(y+1)}.
\end{equation}
Thus $\mu(y)=L/(y+1)+O(y^{-4})$. If $\mu$ were Stieltjes,
Lemma~\ref{lem:rigidity} would force equality with $L/(y+1)$,
contrary to the strict lower inequality in
\eqref{eq:strict-asymptotic}.

For $\mu^2$, the estimate $0<\mu<L/y$ gives
$y\mu(y)^2\to0$. A nonzero Stieltjes function tending to zero
cannot have this behavior: monotone convergence gives
$\lim_{y\to\infty}yV(y)=\rho([0,\infty))>0$, possibly infinite.
This also proves the already-known negative Stieltjes conclusion
for the square without an additional external result.
\end{proof}

\section{A consequence for Robin eigenvalues of boxes}

For a fixed $\beta>0$, consider the first eigenvalue of
\[
 -\Delta u=\lambda_1 u\ \text{in }\Omega,\qquad
 \beta\partial_\nu u+u=0\ \text{on }\partial\Omega,
\]
with the outward normal $\nu$. Keady and Wiwatanapataphee
\cite[Section~6, pp.~924--925]{KW2018} ask about complete
monotonicity under dilation, including the rectangle case.

\begin{corollary}[Dilation of a box]\label{cor:Robin}
For every integer $N\ge1$, every box
$\Omega=\prod_{j=1}^N(-a_j,a_j)$ with $a_j>0$, and fixed
$\beta>0$, the function
$s\mapsto\lambda_1(s\Omega;\beta)$ is strictly completely
monotone on $(0,\infty)$.
\end{corollary}

\begin{proof}
The positive first eigenfunction on an interval $(-a,a)$ is
$\cos(kx)$ with $0<ka<\pi/2$ and
$\beta k\tan(ka)=1$. Equivalently,
\[
 \frac a\beta(\beta k)=\arctan\frac1{\beta k},
 \qquad \beta k=\mu(a/\beta).
\]
Separation of variables for the box gives
\begin{equation}\label{eq:Robin-box}
 \lambda_1(s\Omega;\beta)
 =\beta^{-2}\sum_{j=1}^N\mu(sa_j/\beta)^2.
\end{equation}
Theorem~\ref{thm:main}, followed by positive rescaling of the
argument and finite summation, gives every required strict sign.
\end{proof}

This corollary concerns boxes in arbitrary dimension. It makes no
claim about the dilation question for arbitrary domains.

\section{Relation to previous work and the corpus}\label{sec:literature}

\paragraph{The original problem.}
Keady, Khajohnsaksumeth, and Wiwatanapataphee explicitly pose
Question 1 in \cite[p.~2]{KKW2018}; their Section~5.2 discusses the
Stieltjes possibility for $\mu$. Their Robin eigenvalue work
\cite{KW2017} supplies the analytic setting. The theorem above
keeps the two original inverse functions and their full positive
domains. The family parameter is an auxiliary proof device, and
does not replace the target with a modified problem.

\paragraph{The closest reciprocal-arctangent result.}
Jovanovi\'c and Treml \cite{JT2021} prove logarithmic complete
monotonicity of $1/\arctan x$, which is the $\lambda=0$ member of
\eqref{eq:F-U}. They also show that this reciprocal is not
Stieltjes. Their work already uses contour integral methods and
nonnegative trigonometric Laplace kernels. Accordingly, neither
contour inversion nor positivity of $1-\cos$ is claimed here as
a historically new principle. The calculation needed here is the
parameter-dependent density \eqref{eq:D}, its exact mass
\eqref{eq:mass}, and the sharp classification reaching
$\lambda=1$. The functional inverse $\mu$ is a different
function from the reciprocal $1/\arctan$.

\paragraph{The inverse-ODE step.}
The implication $h'=U(h)$ with $U$ completely monotone implies
$h$ Bernstein is established in the literature; see
\cite[Lemma~2.6]{Li2026} and \cite[Lemma~2.6]{SalazarBF2026}.
We have proved it again for completeness,
not presented it as a new theorem. It supplies a reusable
transfer once the specific kernel positivity has been established.

\paragraph{Connection to the supplied corpus.}
The public corpus \cite{VMRC}, at commit
\texttt{8219b3eb00b8} (the full identifier is in the reference), contains 95
TeX manuscripts. The pertinent methodological comparison is with
manuscript [029], \emph{Observation and Recovery (Geometry,
Mixtures, and Spectral Inversion)}: its parts on inverse moments,
Thorin measures, and mass constraints combine boundary inversion
with exact mass identities to decide positivity. Here those ideas
are used in a squared spectral variable and then transferred to
a functional inverse. No unverified corpus theorem is a premise
of the proof above. All required representations, endpoint
estimates, and sign inductions have been proved in this manuscript.

\paragraph{Scope of the literature check.}
A targeted search through 31 October 2026 examined the original
question, its Robin-eigenvalue context, the reciprocal-arctangent
paper, the supplied corpus, the later papers
\cite{SalazarCM2026,SalazarBF2026}, and the relevant CM/Stieltjes/
Bernstein research archive. The archive explicitly retains the
special Robin inverse branches as open in its dedicated frontier
entry \cite{PudimFrontier}; its negative answer to the different
general inverse-closure Question 3 is not an answer to Question 1.
No identical solution to the specific Question 1 or the sharp
$U_\lambda$ classification was located in the material checked.
This is a bounded literature-search finding, not a proof that no
unindexed or unpublished antecedent exists. The present document
has not undergone external peer review.

\appendix
\section{A second exact representation at the critical endpoint}

The endpoint $\lambda=1$ admits a direct Stieltjes representation
for a different auxiliary function. It offers an independent way
to check the sign and normalization of the central calculation.
Set
\[
 S(w)=\int_0^1\frac{\dd t}{1+wt^2}+\frac1{1+w},\qquad
 B(w)=\frac1{S(w)}-\frac{\sqrt w}{L}.
\]
In the upper half-plane $\operatorname{Im}S(w)<0$; conjugation
and $S(w)>0$ on the positive real axis show that $S$ is
zero-free throughout the slit plane. At zero,
$B(w)=1/2-\sqrt w/L+O(w)$. At $w=-1$ it
is bounded, because $S(w)=(1+w)^{-1}+O(\log(1+w))$. Uniformly
at infinity in the slit plane,
\[
 B(w)=\frac{2}{3L^2w}+O(|w|^{-2}).
\]
For $s>0$, its upper boundary value has imaginary part
$-\pi\rho(s)$, where
\begin{equation}\label{eq:rho}
 \rho(s)=\frac{\sqrt s}{\pi}D_1(\sqrt s)
 =\begin{cases}
 \displaystyle\frac{\sqrt s}{\pi L},&0<s\le1,\\[4pt]
 \displaystyle\frac{\sqrt s\,b_1(\sqrt s)^2}
 {\pi L\bigl(L^2+b_1(\sqrt s)^2\bigr)},&s>1.
 \end{cases}
\end{equation}
The same keyhole-contour argument as in Lemma~\ref{lem:K} gives
\[
 B(w)=\int_0^\infty\frac{\rho(s)}{w+s}\dd s.
\]
Here $\rho(s)=O(\sqrt s)$ near zero and $O(s^{-5/2})$ at
infinity. There is no atom at $s=1$, because $B$ is bounded
there. Since $B(x^2)=x(U_1(x)-1/L)$,
\[
 U_1(x)=\frac1L+\int_0^\infty
                 \frac{\rho(s)}{x(x^2+s)}\dd s.
\]
Using
$[x(x^2+s)]^{-1}=s^{-1}\int_0^\infty e^{-xt}
(1-\cos(\sqrt s\,t))\dd t$ recovers \eqref{eq:endpoint}.
Moreover, $B(0)=1/2$ gives
$\int_0^\infty\rho(s)/s\dd s=1/2$, exactly equivalent to
$\int_0^\infty D_1(r)\dd r=\pi/4$.

\section{Proof obligations and provenance}

\begin{center}
\small
\begin{tabular}{p{0.30\textwidth}p{0.62\textwidth}}
\toprule
Requirement & What is established here \\
\midrule
Pre-existing identified problem & Question 1 of the 2018 article, with its DOI, page, functions, and domain. \\
Original problem unchanged & Both original CM questions and the Stieltjes clause are answered. \\
Previously open status & Explicit in the 2018 paper; a separate 2026 archive still marks the special branches open. Exhaustiveness of the literature search is not claimed. \\
Unconditional full proof & Explicit densities, convergent integrals, exact mass, all-order induction, and the non-Stieltjes proof. \\
More than a restatement & A sharp parameter classification and its endpoint application; the known contour tools and inverse-ODE criterion are credited. \\
No identical prior solution & None located in the documented search. Absolute publication-priority certification remains outside what a finite search can prove. \\
\bottomrule
\end{tabular}
\end{center}

The proof uses no finite-order numerical test as a premise.
In particular, the passage from a positive Laplace density to all
derivative orders follows from a bounded density and dominated
differentiation, while the inverse step is an induction for every
integer order. The source is one standalone TeX file with an inline
bibliography and no external figures or data dependencies.
This manuscript was prepared in an AI-assisted research session;
the explicit arguments are supplied for independent mathematical
and bibliographic review.

\begin{thebibliography}{99}

\bibitem{KKW2018}
G. Keady, N. Khajohnsaksumeth, and B. Wiwatanapataphee,
\emph{On functions and inverses, both positive, decreasing and
convex: And Stieltjes functions}, Cogent Mathematics \& Statistics
\textbf{5} (2018), article 1477543.
\href{https://doi.org/10.1080/25742558.2018.1477543}{doi:10.1080/25742558.2018.1477543}.

\bibitem{KW2017}
G. Keady and B. Wiwatanapataphee,
\emph{Inequalities for the fundamental Robin eigenvalue of the
Laplacian for box-shaped domains}, preprint (2017).
\href{https://arxiv.org/abs/1705.09147}{arXiv:1705.09147}.

\bibitem{KW2018}
G. Keady and B. Wiwatanapataphee,
\emph{Inequalities for the fundamental Robin eigenvalue for the
Laplacian on $N$-dimensional rectangular parallelepipeds},
Mathematical Inequalities \& Applications \textbf{21}(4)
(2018), 911--930.
\href{https://doi.org/10.7153/mia-2018-21-62}{doi:10.7153/mia-2018-21-62}.

\bibitem{JT2021}
V. Jovanovi\'c and M. Treml,
\emph{Logarithmically complete monotonicity of reciprocal arctan
function}, Kragujevac Journal of Mathematics \textbf{49}(1)
(2025), 105--110; preprint (2021).
\href{https://arxiv.org/abs/2112.09960}{arXiv:2112.09960};
\href{https://doi.org/10.46793/KgJMat2501.105J}{doi:10.46793/KgJMat2501.105J}.

\bibitem{Li2026}
C.-G. Li,
\emph{Space-Time Duality of Diffusion via Subordination: General
Theory and Relativistic Models}, preprint, version 2 (2026).
\href{https://arxiv.org/abs/2606.04270}{arXiv:2606.04270}.

\bibitem{SalazarCM2026}
D. S. P. Salazar,
\emph{Determinacy Witnesses in the Completely
Monotone--Stieltjes--Bernstein Hierarchy}, preprint, version~1 (2026).
\href{https://arxiv.org/abs/2607.09868}{arXiv:2607.09868}.

\bibitem{SalazarBF2026}
D. S. P. Salazar,
\emph{Bernstein Functions at Work: Coalescents, Copulas, and
Subordination}, preprint, version~1 (2026).
\href{https://arxiv.org/abs/2607.04467}{arXiv:2607.04467}.

\bibitem{SSV2012}
R. L. Schilling, R. Song, and Z. Vondra\v{c}ek,
\emph{Bernstein Functions: Theory and Applications}, second
edition, De Gruyter, 2012.
\href{https://doi.org/10.1515/9783110269338}{doi:10.1515/9783110269338}.

\bibitem{Pudim2026}
Pudim AI Project,
\emph{Zeta-Law Entropy, Gamma Curvature, and Laplace-Measure
Applications}, public theory archive, latest staged theory v036
(2026), especially APP-0035 and APP-0066. Retrieved 31 October 2026.
\href{https://github.com/pudim-project/pudim-ai-demo-zetalaw}{Public repository}.

\bibitem{PudimFrontier}
Pudim AI Project,
\emph{Keady Robin special inverse branch complete monotonicity
remains open}, public frontier entry, repository commit
\texttt{0867904e9718c33537197541da62c383b8734321}.
\href{https://github.com/pudim-project/pudim-ai-demo-zetalaw/blob/0867904e9718c33537197541da62c383b8734321/wiki/latest/wiki/nodes/T-Keady-special-Robin-inverse-CM-frontier.md}{Versioned frontier entry}.
Retrieved 31 October 2026.

\bibitem{VMRC}
QEEDDD,
\emph{Verifiable Mathematical Reasoning Corpus}, public repository,
commit \texttt{8219b3eb00b83a9bb81bc7b7afbc3c49b4538c9c}
(29 October 2026). In particular manuscript [029],
\emph{Observation and Recovery (Geometry, Mixtures, and Spectral
Inversion)}.
\href{https://github.com/QEEDDD/Verifiable-Mathematical-Reasoning-Corpus}{Repository}.

\end{thebibliography}
\end{document}
